% GENERATED FILE. Edit canonical lessons, then run npm run generate:books.
% canonical-source-hash: fnv1a64:c68526f7b059883f
% canonical-lessons: KA-C460-vasati, KA-C460-ladj, KA-C460-kothadi, KA-C460-svagatakara, KA-C460-kauntar

\chapter{Finding a Room}
\label{ch:ka-finding-a-room}
\hlchaptermodality{460}

\begin{chapteropening}
\textbf{By the end of this chapter:} \emph{I can say lodging, a lodge, a chamber, a receptionist and a counter.}

Complete the last lesson of chapter 460: I can say lodging, a lodge, a chamber, a receptionist and a counter.
\end{chapteropening}

\section[\texorpdfstring{\kn{ವಸತಿ} (vasati)}{ವಸತಿ (vasati)}]{\kn{ವಸತಿ} (vasati) — lodging; accommodation}

\label{lesson:KA-C460-vasati}



\begin{quote}
Before the new one: say the Kannada for a tent, then the Kannada for an encampment.
\end{quote}

\subsection*{You'll want to know: \kn{ವಸತಿ}}
\textbf{\kn{ವಸತಿ}} — \emph{vasati} — \textquotedblleft{}lodging; accommodation\textquotedblright{}.

\begin{quote}
\kn{ಹಬ್ಬದ} \kn{ಸಮಯದಲ್ಲಿ} \kn{ವಸತಿ} \kn{ಸಿಗುವುದು} \kn{ಕಷ್ಟ}.
\end{quote}

\begin{quote}
\emph{habbada samayadalli vasati siguvudu kaṣṭa.}
\end{quote}

\begin{quote}
At festival time lodging is hard to find.
\end{quote}

\kn{ವಸತಿ} \kn{ಶಾಲೆ}, a boarding school; \kn{ವಸತಿಗೃಹ}, a lodge.

\begin{cousinweb}
\textbf{\kn{ವಸತಿ}} (\emph{vasati}) is borrowed from Sanskrit \emph{vasati}, \textquotedblleft{}a dwelling, a lodging\textquotedblright{}, from \emph{vas}, \textquotedblleft{}to stay, to dwell\textquotedblright{} (CDIAL 11435).
\end{cousinweb}

\begin{grammarlens}[title={What you've built}]
One more word for getting around.
\end{grammarlens}

\subsection*{Your turn}
\begin{itemize}
\raggedright
  \item \emph{Say it:} \emph{vasati}
  \item \emph{Say it:} \emph{vasati}, once more
  \item \emph{Say it:} \emph{habbada samayadalli vasati siguvudu kaṣṭa.} — At festival time lodging is hard to find.
  \item \emph{Recall:} say the Kannada for a tent, then the Kannada for an encampment, then say \emph{vasati} again
\end{itemize}

\begin{tcolorbox}[breakable,colback=teal!4,colframe=teal!35!black,title={Before you move on}]
What does \kn{ವಸತಿ} mean? (\textquotedblleft{}Lodging; accommodation\textquotedblright{}.) Say it once more.
\end{tcolorbox}
\section[\texorpdfstring{\kn{ಲಾಡ್ಜ್} (lāḍj)}{ಲಾಡ್ಜ್ (lāḍj)}]{\kn{ಲಾಡ್ಜ್} (lāḍj) — a lodge; a cheap hotel}

\label{lesson:KA-C460-ladj}



\begin{quote}
Before the new one: say the Kannada for an encampment, then the Kannada for lodging.
\end{quote}

\subsection*{You'll want to know: \kn{ಲಾಡ್ಜ್}}
\textbf{\kn{ಲಾಡ್ಜ್}} — \emph{lāḍj} — \textquotedblleft{}a lodge; a cheap hotel\textquotedblright{}.

\begin{quote}
\kn{ಬಸ್} \kn{ನಿಲ್ದಾಣದ} \kn{ಹತ್ತಿರ} \kn{ಒಂದು} \kn{ಲಾಡ್ಜ್} \kn{ಇದೆ}.
\end{quote}

\begin{quote}
\emph{bas nildāṇada hattira ondu lāḍj ide.}
\end{quote}

\begin{quote}
There is a lodge near the bus station.
\end{quote}

A plain place to sleep; a \kn{ಹೋಟೆಲ್} may also serve food.

\begin{cousinweb}
\textbf{\kn{ಲಾಡ್ಜ್}} (\emph{lāḍj}) is a loan from English \emph{lodge}.
\end{cousinweb}

\begin{grammarlens}[title={What you've built}]
One more word for getting around.
\end{grammarlens}

\subsection*{Your turn}
\begin{itemize}
\raggedright
  \item \emph{Say it:} \emph{lāḍj}
  \item \emph{Say it:} \emph{lāḍj}, once more
  \item \emph{Say it:} \emph{bas nildāṇada hattira ondu lāḍj ide.} — There is a lodge near the bus station.
  \item \emph{Recall:} say the Kannada for an encampment, then the Kannada for lodging, then say \emph{lāḍj} again
\end{itemize}

\begin{tcolorbox}[breakable,colback=teal!4,colframe=teal!35!black,title={Before you move on}]
What does \kn{ಲಾಡ್ಜ್} mean? (\textquotedblleft{}A lodge; a cheap hotel\textquotedblright{}.) Say it once more.
\end{tcolorbox}
\section[\texorpdfstring{\kn{ಕೊಠಡಿ} (koṭhaḍi)}{ಕೊಠಡಿ (koṭhaḍi)}]{\kn{ಕೊಠಡಿ} (koṭhaḍi) — a chamber; a room}

\label{lesson:KA-C460-kothadi}



\begin{quote}
Before the new one: say the Kannada for lodging, then the Kannada for a lodge.
\end{quote}

\subsection*{You'll want to know: \kn{ಕೊಠಡಿ}}
\textbf{\kn{ಕೊಠಡಿ}} — \emph{koṭhaḍi} — \textquotedblleft{}a chamber; a room\textquotedblright{}.

\begin{quote}
\kn{ನಮ್ಮ} \kn{ಕೊಠಡಿ} \kn{ಮೂರನೇ} \kn{ಮಹಡಿಯಲ್ಲಿ} \kn{ಇತ್ತು}.
\end{quote}

\begin{quote}
\emph{namma koṭhaḍi mūranē mahaḍiyalli ittu.}
\end{quote}

\begin{quote}
Our room was on the third floor.
\end{quote}

A formal word; \kn{ಕೋಣೆ} is the everyday room.

\begin{cousinweb}
\textbf{\kn{ಕೊಠಡಿ}} (\emph{koṭhaḍi}) is usually traced to Sanskrit \emph{kōṣṭha}, \textquotedblleft{}a storeroom, an inner apartment\textquotedblright{} (CDIAL 3546), through Prakrit.
\end{cousinweb}

\begin{grammarlens}[title={What you've built}]
One more word for getting around.
\end{grammarlens}

\subsection*{Your turn}
\begin{itemize}
\raggedright
  \item \emph{Say it:} \emph{koṭhaḍi}
  \item \emph{Say it:} \emph{koṭhaḍi}, once more
  \item \emph{Say it:} \emph{namma koṭhaḍi mūranē mahaḍiyalli ittu.} — Our room was on the third floor.
  \item \emph{Recall:} say the Kannada for lodging, then the Kannada for a lodge, then say \emph{koṭhaḍi} again
\end{itemize}

\begin{tcolorbox}[breakable,colback=teal!4,colframe=teal!35!black,title={Before you move on}]
What does \kn{ಕೊಠಡಿ} mean? (\textquotedblleft{}A chamber; a room\textquotedblright{}.) Say it once more.
\end{tcolorbox}
\section[\texorpdfstring{\kn{ಸ್ವಾಗತಕಾರ} (svāgatakāra)}{ಸ್ವಾಗತಕಾರ (svāgatakāra)}]{\kn{ಸ್ವಾಗತಕಾರ} (svāgatakāra) — a receptionist}

\label{lesson:KA-C460-svagatakara}



\begin{quote}
Before the new one: say the Kannada for a lodge, then the Kannada for a chamber.
\end{quote}

\subsection*{You'll want to know: \kn{ಸ್ವಾಗತಕಾರ}}
\textbf{\kn{ಸ್ವಾಗತಕಾರ}} — \emph{svāgatakāra} — \textquotedblleft{}a receptionist\textquotedblright{}.

\begin{quote}
\kn{ಸ್ವಾಗತಕಾರ} \kn{ನಮಗೆ} \kn{ಕೊಠಡಿಯ} \kn{ಕೀಲಿ} \kn{ಕೊಟ್ಟನು}.
\end{quote}

\begin{quote}
\emph{svāgatakāra namage koṭhaḍiya kīli koṭṭanu.}
\end{quote}

\begin{quote}
The receptionist gave us the key to the room.
\end{quote}

\kn{ಸ್ವಾಗತ}, welcome, plus -\kn{ಕಾರ}, one who does it.

\begin{cousinweb}
\textbf{\kn{ಸ್ವಾಗತಕಾರ}} (\emph{svāgatakāra}) is built from Sanskrit \emph{svāgata}, \textquotedblleft{}welcome\textquotedblright{} (CDIAL 13920), and \emph{-kāra}, \textquotedblleft{}one who does\textquotedblright{}: the one who welcomes.
\end{cousinweb}

\begin{grammarlens}[title={What you've built}]
One more word for getting around.
\end{grammarlens}

\subsection*{Your turn}
\begin{itemize}
\raggedright
  \item \emph{Say it:} \emph{svāgatakāra}
  \item \emph{Say it:} \emph{svāgatakāra}, once more
  \item \emph{Say it:} \emph{svāgatakāra namage koṭhaḍiya kīli koṭṭanu.} — The receptionist gave us the key to the room.
  \item \emph{Recall:} say the Kannada for a lodge, then the Kannada for a chamber, then say \emph{svāgatakāra} again
\end{itemize}

\begin{tcolorbox}[breakable,colback=teal!4,colframe=teal!35!black,title={Before you move on}]
What does \kn{ಸ್ವಾಗತಕಾರ} mean? (\textquotedblleft{}A receptionist\textquotedblright{}.) Say it once more.
\end{tcolorbox}
\section[\texorpdfstring{\kn{ಕೌಂಟರ್} (kauṇṭar)}{ಕೌಂಟರ್ (kauṇṭar)}]{\kn{ಕೌಂಟರ್} (kauṇṭar) — a counter; a desk}

\label{lesson:KA-C460-kauntar}



\begin{quote}
Before the new one: say the Kannada for a chamber, then the Kannada for a receptionist.
\end{quote}

\subsection*{You'll want to know: \kn{ಕೌಂಟರ್}}
\textbf{\kn{ಕೌಂಟರ್}} — \emph{kauṇṭar} — \textquotedblleft{}a counter; a desk\textquotedblright{}.

\begin{quote}
\kn{ಟಿಕೆಟ್} \kn{ಕೌಂಟರ್ನಲ್ಲಿ} \kn{ಉದ್ದ} \kn{ಸಾಲು} \kn{ಇತ್ತು}.
\end{quote}

\begin{quote}
\emph{ṭikeṭ kauṇṭarnalli udda sālu ittu.}
\end{quote}

\begin{quote}
There was a long queue at the ticket counter.
\end{quote}

\kn{ಸ್ವಾಗತ} \kn{ಕೌಂಟರ್}, a reception desk.

\begin{cousinweb}
\textbf{\kn{ಕೌಂಟರ್}} (\emph{kauṇṭar}) is a loan from English \emph{counter}.
\end{cousinweb}

\begin{grammarlens}[title={What you've built}]
One more word for getting around.
\end{grammarlens}

\subsection*{Your turn}
\begin{itemize}
\raggedright
  \item \emph{Say it:} \emph{kauṇṭar}
  \item \emph{Say it:} \emph{kauṇṭar}, once more
  \item \emph{Say it:} \emph{ṭikeṭ kauṇṭarnalli udda sālu ittu.} — There was a long queue at the ticket counter.
  \item \emph{Recall:} say the Kannada for a chamber, then the Kannada for a receptionist, then say \emph{kauṇṭar} again
\end{itemize}

\begin{tcolorbox}[breakable,colback=teal!4,colframe=teal!35!black,title={Before you move on}]
What does \kn{ಕೌಂಟರ್} mean? (\textquotedblleft{}A counter; a desk\textquotedblright{}.) Say it once more.
\end{tcolorbox}
